We performed a post hoc sensitivity analysis of the frozen common-candidate allocation study, after its original test results were known. Within each original best-validation-BER tie set (tolerance $10^{-12}$), choosing the smallest or largest relation coefficient changed 86 and 100 of 300 selections, respectively. The original default-first results remain primary. Across the 20 independent-selection ACC contrasts, Full had a positive mean difference in 18 and 16 contrasts under the smallest- and largest-coefficient rules, compared with 19 under the original rule. The corresponding counts for controls using each policy's Full-selected coefficient were 19 and 19, compared with 19 originally. All three policies retain the same data, query streams, eligible folds and coefficient grid. This analysis describes dependence on tie breaking within existing corpora; it supplies neither independent confirmation nor a test-based choice of a preferred policy.
The exact 19-of-20 positive-mean count depended on the tie rule. A majority of the independent-selection ACC mean contrasts remained positive under all three rules. These sign counts do not establish uniform improvement across runs, metrics, or untested tie rules.
